\section{Packet scores and finite decompositions}
\label{sec:packet-structure}

The classical theorem controls positive masses transported along exact
lines. For oscillatory packets, an observed mass is the square of a
coherent sum, and an earlier packet reaches the output within its
own position uncertainty, up to controlled tails. We therefore begin with analytic mass and
decomposition estimates before introducing any positive auxiliary law.
The packet score below is designed to compose across a time cut.
This section establishes its preliminary growth bound; the next two
sections reduce excess growth to the uncertainty boundary and exclude
it by a finite information argument.

The transverse dimension is two and $0<\gamma<1$.
All entropies are normalized by $\ell=\log(1/\eps)$, and
$\logeps r=\log r/\ell$. Constants in analytic estimates are
independent of $\eps$ and of the input function.

\subsection{The phase class and the growth supremum}

The longitudinal coordinate is \(c\in[0,1]\), the input variable is
\(v\in\RR^2\), and the observation position is \(x\in\RR^2\).
The phase is defined on \([0,1]\times\RR^2\times U_v\), for an open
neighborhood \(U_v\) of the bounded input patch.  We use real phases satisfying
\begin{equation}
 \partial_v\Phi(c,x,v)=\zeta(v,x-cv).
 \label{eq:packet-phase}
\end{equation}
The function \(\zeta\) is defined for every \(p\in\RR^2\) and on all
bounded \(v\)-patches needed in the argument.  On each such patch,
\(\partial_p\zeta(v,p)\) is symmetric and either
\begin{equation}
 a_0 I\le \partial_p\zeta(v,p)\le a_1 I
 \quad\hbox{for all }(v,p),
 \label{eq:packet-definiteness}
\end{equation}
or the negative of this inequality holds.  Here \(a_0>0\).
For some fixed \(A_\Phi\ge1\), all derivatives of order \(j\) of this
matrix are bounded by \(A_\Phi^{j+1}(j!)^3\).  Bounds on domains,
\(a_0^{-1}\), \(a_1\), and \(A_\Phi\) may vary between applications;
they are uniform within each sequence in which \(\eps\to0\).
Input-only terms in \(\Phi\) are absorbed into the input.

Write
\begin{equation}
 \lambda^{-1}=\eps^{\Planck},\qquad
 u_t=\frac{\Planck-t}{2},\qquad
 \rho_t=\eps^{u_t},\qquad
 \sigma_t=(\lambda\rho_t)^{-1}=\eps^{u_t+t}.
 \label{eq:packet-scales}
\end{equation}
Time boxes are nested dyadic intervals in \([0,1]\), with root center
\(C_0=1/2\).  A depth-\(t\) box has length comparable to \(\eps^t\).
Its center is denoted by \(C_t\).  We take actual dyadic lengths in
changes of coordinates; this changes depths by at most
\(\log 2/\ell\).

An observation \(O=(C_t,V,X)\) uses velocity and position lattices of
spacings \(\rho_t\) and \(\sigma_t\), respectively.  Bounded
multiplicities and translated lattices are allowed.  Its mass is
\begin{equation}
 m_O(f)=\left|\rho_t^{-1}\int f(v)
  \psi_O\left(\frac{v-V}{\rho_t}\right)
  e^{i\lambda\Phi(C_t,X,v)}\,dv\right|^2.
 \label{eq:packet-mass}
\end{equation}
The input belongs to \(L^2\) and has bounded support.  The amplitudes
have fixed compact supports in their displayed variables, and every
derivative has a uniform bound.  The centers \(X,V\) range over bounded
sets when \(t>0\).  Suprema over any specified amplitude class with
these uniform bounds are also permitted.  Lemma~\ref{lem:packet-bessel}
justifies these suprema by simultaneous choices at the sites.

At the root, the canonical mass is the supremum in
\eqref{eq:packet-mass} over amplitudes supported in \(\{|z|\le100\}\)
with
\(\|\partial^\alpha\psi\|_\infty
 \le100^{|\alpha|+1}(|\alpha|!)^{10}\).
The root grids are fixed.  The position cutoff is
\(|X|\le1+\log\lambda\), with a natural logarithm here, and velocity
centers lie within \(100\rho_0\) of the input support, or in a fixed
sufficiently large enlargement of that support.

We use the aperture convention of
Definition~\ref{def:classical-aperture}.  Thus a label \(D\) includes
\(C_t\), an orientation, a rectangle locating \(V\), and a rectangle
locating \(X\), with tangent/normal widths
\((\eps^b,\eps^a)\) and
\((\eps^{b+t},\eps^{a+t})\).  Labels may contain additional information.
For packets the admissible depths satisfy
\begin{equation}
 0\le b\le a\le\min(L-t,u_t),\qquad L\ge t>0,
 \label{eq:packet-apertures}
\end{equation}
and \(w(b,a)=(1+\gamma)b+(1-\gamma)a\).
For any probability law on observations and aperture labels, supported
where \(m_O>0\), define
\begin{equation}
 P_t=\Ent(O)+\frac12\Ent(O\mid D)
       +\frac32\EE\logeps{m_O}
       +\frac12\EE w(D).
 \label{eq:packet-score}
\end{equation}
Let \(K_0(f)\) be the supremum of the corresponding root scores for
the canonical masses, with \(0\le b\le a\le\min(L,u_0)\),
enlarged to include the value
\(\frac32\logeps{\|f\|_2^2}\).  The last option is part of the
definition, including for rescaled local inputs.

The entropy terms connect the score to the desired \(L^3\) estimate.
For a finite family with at least one positive packet mass, omit the
zero masses. The log-sum identity gives
\[
 \sup_{\pi}\left\{\Ent_\pi(O)
       +\frac32\EE_\pi\logeps{m_O}\right\}
 =\logeps{\sum_Om_O^{3/2}},
\]
where the supremum is over probability laws on the remaining sites.
Because \(m_O\) is a squared packet amplitude, the right side is the
normalized logarithm of the cubic packet sum. At \(t=\Planck\),
admissibility forces \(b=a=0\), hence \(w=0\). For the maximizing
observation law, every compatible endpoint score is at least the
displayed value, because \(\Ent(O\mid D)\ge0\).
For a given aperture score, put \(M=\sum_Om_O\), and let \(M_D\) be the
sum over all sites in the cap specified by \(D\). Applying log-sum once
without conditioning and once conditional on \(D\) gives
\[
 P_t\le \logeps{M}
       +\frac12\EE\bigl[\logeps{M_D}+w(D)\bigr].
\]
Thus total and cap masses control the root aperture scores; the
additional norm option in \(K_0\) is estimated separately.

Multiplying \(f\) by a scalar shifts both \(P_t\) and \(K_0\) by the
same quantity.  We may therefore normalize \(\|f\|_2=1\) when considering
their difference.  Define \(\betaq\) to be the supremum of
\begin{equation}
 \limsup_{\eps\to0}\frac{P_t-K_0}{t}
 \label{eq:packet-growth-supremum}
\end{equation}
over all the preceding systems, allowing convergent bounded depth
parameters with positive limiting duration.  In each sequence all
domain, phase, amplitude, and multiplicity bounds are fixed.  The
input and the individual phase and amplitude choices may depend on
\(\eps\).  The supremum over sequences also allows any choices of
these fixed bounds.

\begin{theorem}[Packet growth]
\label{thm:packet-growth}
For the phase class \eqref{eq:packet-phase}--\eqref{eq:packet-definiteness},
\(\betaq\le\gamma/2\).  Equivalently, for fixed admissible depth
parameters and uniform bounds on the data, and for every \(\nu>0\),
there is \(\eps_0>0\) such that
\begin{equation}
 P_t\le K_0+\left(\frac\gamma2+\nu\right)t
 \qquad(0<\eps<\eps_0).
 \label{eq:packet-growth-finite}
\end{equation}
The same statement holds for convergent bounded parameters with
positive limiting duration.
\end{theorem}

The proof is completed in Section~\ref{sec:packet-repayment}.
The present section proves all the analytic decomposition statements
and the preliminary bound \(\betaq\le1\) used in that proof.
Uniformity in \eqref{eq:packet-growth-finite} follows from the supremum
definition: a failure for fixed bounds would select a sequence
contradicting the claimed bound on \(\betaq\).

\subsection{Bessel estimates and changes of scale}

Fourier localization and almost orthogonality are the standard analytic
ingredients of wave-packet decomposition; see, for example,
\citet{Tao2003} and
\citet{GuoOhWangWuZhang2025}. We give the estimates
with the phase class and uniform amplitude bounds required by our
packet score, including changes of scale and quantified tails.

\begin{lemma}[Uniform packet Bessel estimates]
\label{lem:packet-bessel}
At a fixed time anchor, packets with natural spacings \(\rho\) in
velocity and \(\sigma=(\lambda\rho)^{-1}\) in position satisfy
\begin{equation}
 \sum_O m_O(f)\le C\|f\|_2^2,\qquad
 \left\|\sum_T c_T g_T\right\|_2^2\le C\sum_T|c_T|^2,
 \label{eq:packet-bessel}
\end{equation}
where \(g_T\) are the normalized synthesis packets with the opposite
phase sign.  These inequalities hold for subsets and for canonical
supremum masses.  Their constants depend only on ellipticity, finitely
many phase and amplitude derivative bounds, support sizes, and lattice
multiplicity, and not on the number of sites or on the position cutoff.

More precisely, packets at this same time and scale obey, for every
integer \(M\ge1\),
\begin{equation}
 |\langle g_{V,X},g_{V',X'}\rangle|
 \le C_M\ind_{\{|V-V'|\le C\rho\}}
       \left(1+\frac{|X-X'|}{\sigma}\right)^{-M}.
 \label{eq:packet-gram}
\end{equation}
At a local root, where both natural spacings equal \(s\), let \(F\)
be any allowed phase-space cap and let \(F^R\) denote its enlargement
by distance \(Rs\) in each velocity and position coordinate,
\(R\ge2\).  For \(h=\sum_Tc_Tg_T\) and every integer \(M>2\),
canonical analysis gives
\begin{equation}
 \sum_{O\in F}m_O(h)
 \le C\sum_{T\in F^R}|c_T|^2
       +C_M R^{4-2M}\sum_T|c_T|^2.
 \label{eq:packet-cap-bessel}
\end{equation}
Here a fixed enlargement of the support sizes is included in \(F^R\).
Every axis of \(F\) is at least \(s\); consequently \(F^R\) is covered
by at most \(CR^4\) caps of the same widths.
\end{lemma}

\begin{proof}
Only overlapping velocity supports contribute.  Put \(v=V+\rho z\).
For the phase difference at equal time, its gradient in \(z\) is
\[
 \lambda\rho\{\zeta(v,X-cv)-\zeta(v,X'-cv)\}.
\]
The braces equal \(B(v)(X-X')\), where \(B(v)\) is the integral of
\(\partial_p\zeta\) along the segment between the two spatial
arguments.  The sign and lower bound in
\eqref{eq:packet-definiteness} give gradient size at least
\(a_0|X-X'|/\sigma\).  Each higher derivative is bounded by
\(C_j|X-X'|/\sigma\).  Repeated use of
\[
 \frac{\nabla\varphi\cdot\nabla_z}{i|\nabla\varphi|^2}
 e^{i\varphi}=e^{i\varphi}
\]
therefore gives \eqref{eq:packet-gram}: derivatives of its coefficients
are bounded by the corresponding inverse gradient powers.  The
estimate with a factor \(1+|X-X'|/\sigma\) also covers bounded
separation.  Taking \(M>2\), the row and column sums of the Gram
matrix are bounded, since there are only boundedly many overlapping
velocity centers and the position lattice is two-dimensional.
The Schur estimate proves synthesis Bessel; its adjoint proves
analysis Bessel.

For supremum masses, choose independently at every site of an arbitrary
finite subset an amplitude whose squared pairing is within a fixed
factor of that site's supremum.  The same Gram estimate is uniform in
all these choices.  Taking the supremum and then exhausting the sites
proves the assertion.

For \eqref{eq:packet-cap-bessel}, split the synthesis coefficients into
\(F^R\) and its complement.  Analysis and synthesis Bessel applied
successively bound the first part by its coefficient square sum.
For the other part, a nonzero Gram entry either has spatial separation
at least a fixed multiple of \(Rs\), or is zero by disjoint velocity
supports.  Summing \eqref{eq:packet-gram} outside that spatial radius
gives row and column sums at most \(C_MR^{2-M}\).  The resulting
operator norm squared is at most \(C_MR^{4-2M}\).  The inequality
\(|a+b|^2\le2|a|^2+2|b|^2\) proves the claim after changing constants.
This argument is again uniform over choices of canonical amplitudes.
The covering count follows by enlarging each of the four coordinates:
an interval of length at least \(s\), enlarged by \(Rs\), needs at
most \(C R\) original-length intervals.
\end{proof}

Moving a center by a bounded number of its lattice spacings changes
only an admissible amplitude factor, after removal of a scalar phase.
The same is true when the time anchor moves inside its box and \(X\)
is transported by that displacement times \(V\).  Indeed, after
a scalar phase is removed, the first derivative is bounded in packet
units; the higher derivatives of the phase change are bounded by
\(C_j\lambda\rho_t^2\eps^t=C_j\).  This also proves the bounded
ambiguity assertions used with shifted grids below.

Fix \(0\le r<t\), fixed widths \(b,a\), and \(0\le b_*\le b\).
The geometric localization label is
\begin{equation}
 J_0=\left(C_r,V_{b_*},
             (X+(C_r-C_t)V)_{r+b_*}\right).
 \label{eq:packet-localizer}
\end{equation}
Subscripts on vectors denote grid cells.  The endpoint observation
and its aperture each determine this label to bounded ambiguity;
adjoining the bounded neighboring-cell choices makes the determination
exact.  Write \(V_J,X_J\) for its velocity center and its trajectory
position at \(C_r\).  Let
\(s_v=\eps^{b_*}\), \(s_c=\eps^r\), \(s_x=s_cs_v\).  In exact
dyadic coordinates take the actual length of the time box for \(s_c\).
Set
\begin{align}
 v&=V_J+s_vv',&
 c&=C_r+s_c(c'-1/2),\\
 x&=X_J+(c-C_r)V_J+s_xx',&
 p'&=x'-c'v'.
 \label{eq:packet-rescale-coordinates}
\end{align}
Then \(x-cv=X_J-C_rv+s_x(p'+v'/2)\).  Subtract
\(\Phi(C_r,X_J,v)\) and divide the phase by \(s_cs_v^2\).
The new gradient function is
\begin{equation}
 \zeta'(v',p')=
 \frac{\zeta(v,X_J-C_rv+s_x(p'+v'/2))
       -\zeta(v,X_J-C_rv)}{s_x}.
 \label{eq:packet-rescaled-phase}
\end{equation}
In particular,
\(\partial_{p'}\zeta'=\partial_p\zeta(v,X_J-C_rv+s_x(p'+v'/2))\).
Symmetry, definiteness, and uniform derivative bounds persist, with
changed fixed constants.  Differentiating this formula gives all
the required derivative bounds; the affine substitutions have bounded
derivatives, and \(s_v,s_x\le1\).
The new parameters and input are
\begin{equation}
 \Planck'=\Planck-r-2b_*,\quad L'=L-r-b_*,\quad
 f'(v')=s_v f(V_J+s_vv')e^{i\lambda\Phi(C_r,X_J,V_J+s_vv')}.
 \label{eq:packet-rescaled-parameters}
\end{equation}
The input change preserves its \(L^2\) norm and the pairing
\eqref{eq:packet-mass}.  At local endpoint time \(t-r\), the velocity
and position depths are \(u_t-b_*\) and \(u_t+t-r-b_*\).
Also \(\Planck'\ge t-r>0\), and the new aperture widths are
\((b-b_*,a-b_*)\).  Their weight is \(w(b,a)-2b_*\).
Conversely, a local root aperture becomes an aperture at \(r\) after
adding \(b_*\) to both widths.  Thus returning a score to the original
coordinates adds \(b_*\) to its weight contribution.
The local inputs used below are sums of windows with centers in a
fixed enlargement of the velocity cell of \(J_0\).  Their supports
are therefore bounded in \(v'\), uniformly in the localizer.

\subsection{A finite Fourier decomposition with quantified tails}

We state error sizes before taking any limits.  At fixed depths and
bounded domains all the observation lattices, retained atom sets,
and cap menus below have at most \(C\eps^{-Q}\) elements for some
fixed \(Q\).  The logarithmic root cutoff contributes a factor which
can also be absorbed in this bound by increasing \(Q\).  This exponent
does not bound the number of arbitrary extra labels; those are
accounted for separately.

\begin{lemma}[Finite packet decomposition]
\label{lem:packet-decomposition}
Fix the data bounds, admissible depths, \(r<t\), and \(b_*\le b\).
Put
\(\Delta=u_r-b_*\); then \(\Delta\ge(t-r)/2>0\).
For each localization \(J_0\), there is a finite family of atoms
\begin{equation}
 g_T(v)=\rho_r^{-1}\eta_T((v-V_T)/\rho_r)
                     e^{-i\lambda\Phi(C_r,X_T,v)},
 \qquad h_J=\sum_Tc_Tg_T,
 \label{eq:packet-earlier-atoms}
\end{equation}
with \(V_T\) within a fixed enlargement of its velocity cell and
\(|X_T-X_J|\le C\eps^{r+b_*}\).  The amplitudes have uniform
derivative bounds, the sites have bounded multiplicity, and
\(\sum_T|c_T|^2\le C\|f\|_2^2\).  Each coefficient is an analysis
test at time \(r\).  If \(r=0\), then
\begin{equation}
 |c_T|^2\le C_{\rm can}m_T^{\rm can}(f)
 \label{eq:packet-coeff-canonical}
\end{equation}
with a fixed constant.

More quantitatively, choose \(0<\eta<\Delta/4\) and any required
power \(A>0\).  There are an integer integration order \(M\), a
coefficient cutoff exponent \(B\), constants \(C_A\), and
\(\eps_0>0\), depending only on these choices and the fixed data,
such that, for \(0<\eps<\eps_0\):

\begin{enumerate}
\item Replacing \(f\) by \(h_J\) in every endpoint pairing belonging
to \(J_0\) changes it by at most \(C_A\eps^A\|f\|_2\).
Coefficients with \(|c_T|^2<\eps^B\|f\|_2^2\) may be omitted with
the same error.
\item An individual endpoint test interacts, up to such an error,
only with overlapping velocity supports and atoms satisfying
\begin{equation}
 |V-V_T|\le C\rho_t,\qquad
 |X-X_T-(C_t-C_r)V_T|\le C\eps^{-\eta}\sigma_r.
 \label{eq:packet-transport}
\end{equation}
This assertion also holds for any subset of the retained coefficients.
\item If one subsequently sums squared errors over at most
\(C\eps^{-Q}\) tests, the orders can be chosen to give total squared
error at most \(C_A\eps^{2A}\|f\|_2^2\).
\end{enumerate}

For example it suffices to choose
\begin{equation}
 M>4+\frac{A+Q+4}{\min(\eta,\Delta/2)},\qquad
 B>2A+3Q+10.
 \label{eq:packet-tail-orders}
\end{equation}
Only finitely many derivatives, up to an order determined by \(M\),
are used.  All constants are fixed before \(\eps\) tends to zero.
The enlargement parameter \(\eta\) may tend to zero afterwards.
After the finite family and its coefficient cutoff \(B\) have been
chosen, the integration order for the transport estimate in part~2
can be increased independently, without changing that family or \(B\).
\end{lemma}

\begin{proof}
Choose real smooth windows \(\chi_j\), at scale \(\rho_r\), with
bounded overlap and \(\sum_j\chi_j^2=1\) on the velocity region
needed by this localizer.  They can be chosen with support inside
a square of side \(4\rho_r\) and rescaled derivative bounds
\(C^{k+1}(k!)^2\).  To construct them, take translates of a
Gevrey bump positive on a lattice cell and divide by the square root
of the positive, boundedly overlapping sum of their squares.
One may build the bump from \(e^{-1/s}\ind_{\{s>0\}}\): Cauchy's
formula on a complex disk of radius a small multiple of \(s\),
followed by maximizing \(s^{-k}e^{-c/s}\), bounds its derivatives
by \(C^{k+1}(k!)^2\).  Products of translated copies provide the
required compactly supported bump.
The product and composition derivative formulas preserve this
Gevrey order, with changed constants.

Write \(z=(v-V_j)/\rho_r\), and take a square of side \(K\), fixed
and larger than the supports of the rescaled windows.  The function
\[
 H_j(z)=\rho_r f(V_j+\rho_r z)\chi_j(z)
              e^{i\lambda\Phi(C_r,X_J,V_j+\rho_r z)}
\]
has an \(L^2\) Fourier series in the orthonormal basis
\(K^{-1}e^{2\pi i n\cdot z/K}\), \(n\in\mathbb Z^2\).
Denote its coefficients by \(c_{j,n}\).  Parseval and
\(\sum_j\chi_j^2=1\) give
\(\sum_{j,n}|c_{j,n}|^2\le\|f\|_2^2\).
Before any truncation, synthesis uses the exact atoms
\begin{equation}
 (K\rho_r)^{-1}\chi_j(z)
 e^{-i\lambda\Phi(C_r,X_J,v)}e^{2\pi i n\cdot z/K}.
 \label{eq:packet-linear-fourier-atoms}
\end{equation}
Their sum is \(f\) on the needed velocity region in \(L^2\).

For each mode, solve
\begin{equation}
 \zeta(V_j,X_{j,n}-C_rV_j)
 =\zeta(V_j,X_J-C_rV_j)-\frac{2\pi n}{K\lambda\rho_r}.
 \label{eq:packet-gradient-matching}
\end{equation}
For fixed \(V_j\), the map on the left is globally invertible:
its definite derivative gives both the lower Lipschitz bound
\(a_0|p-q|\) and the upper bound \(a_1|p-q|\); it is a local
diffeomorphism and is proper, hence its open and closed image is
all of \(\RR^2\).  It is injective by the lower bound.
The mode-to-position correspondence is therefore bilipschitz at
spacing \(\sigma_r\).  Round positions to that lattice, retaining
the original mode as a bounded multiplicity label.

For retained centers with \(|X_{j,n}-X_J|\le C\eps^{r+b_*}\),
\eqref{eq:packet-linear-fourier-atoms} equals
\eqref{eq:packet-earlier-atoms} with amplitude
\[
 \eta_{j,n}(z)=K^{-1}\chi_j(z)
 \exp\left(i\lambda[\Phi(C_r,X_{j,n},v)-\Phi(C_r,X_J,v)]
                    +\frac{2\pi i n\cdot z}{K}\right).
\]
The linear term cancels at \(V_j\), up to the bounded error caused
by rounding.  For derivatives of order \(k\ge2\), the exponent
has bounds controlled by
\[
 C_k\lambda\rho_r^k\eps^{r+b_*}
 \le C_k\lambda\rho_r^2\eps^r=C_k.
\]
Here and below the derivatives are obtained by subtracting the two
\(\zeta\) expressions and integrating \(\partial_p\zeta\);
this also bounds mixed derivatives involving \(v\).
We used the exact scale identity
\(\lambda\rho_r^2\eps^r=1\), up to the fixed dyadic rounding
factor.  Thus the amplitude bounds are uniform.
The same Fourier coefficient, expressed with the conjugate of this
amplitude, is an analysis pairing at \((C_r,V_j,X_{j,n})\).

The derivative bounds on the exponent are, more precisely,
\(C^{k+1}(k!)^3\); the product and exponential derivative formulas
give the deliberately weaker but convenient bound
\(C_1^{k+1}(k!)^5\) on \(\eta_{j,n}\).
Since
\[
 \sup_{k\ge0}\frac{C_1^{k+1}}{100^{k+1}(k!)^5}<\infty,
\]
a fixed scalar multiple of every such amplitude belongs to the
canonical class.  Its support is inside the radius-100 ball.
At \(r=0\) the retained spatial centers are bounded and hence lie
inside \(1+\log\lambda\) for small \(\eps\).
This proves \eqref{eq:packet-coeff-canonical}.

We now estimate the discarded modes in their linear form
\eqref{eq:packet-linear-fourier-atoms}, before replacing their phase
by a packet phase.  On a window meeting an endpoint velocity support,
the phase difference between \((C_t,X)\) and \((C_r,X_J)\)
has all derivatives of order at least two in \(z\) bounded by
\(C_k\).  Indeed the spatial arguments differ by
\(O(\eps^{r+b_*})\), and the time difference is \(O(\eps^r)\);
the factor \(\lambda\rho_r^2\) is canceled by \(\eps^r\).
By \eqref{eq:packet-gradient-matching} and definiteness, the gradient
at the window center has size comparable, up to bounded errors, to
\[
 d_{j,n,O}=
 \frac{|X-X_{j,n}-(C_t-C_r)V_j|}{\sigma_r}.
\]
Its higher derivatives have bounds independent of \(n\).  The same
integration operator as in Lemma~\ref{lem:packet-bessel} therefore
gives, for every \(M\),
\begin{equation}
 |\langle g_{j,n},h_O\rangle|
 \le C_M\frac{\rho_r}{\rho_t}(1+d_{j,n,O})^{-M},
 \label{eq:packet-mixed-gram}
\end{equation}
where \(h_O\) denotes the normalized test.  This notation refers
to the exact linear atom when its center has not yet been retained.
The factor \(\rho_r/\rho_t\) is the product of the two normalization
factors and the window area.  Derivatives of the endpoint amplitude
in window units are uniformly bounded because \(\rho_r\le\rho_t\).

For each window, the bilipschitz mode lattice gives
\[
 \sum_{n:d_{j,n,O}>R}|\langle g_{j,n},h_O\rangle|^2
 \le C_M(\rho_r/\rho_t)^2 R^{2-2M}\qquad(R\ge2).
\]
At most \(C(\rho_t/\rho_r)^2\) windows meet an endpoint test.
Summing cancels these scale factors.  Cauchy--Schwarz with Parseval
then bounds the discarded pairing by
\begin{equation}
 C_MR^{1-M}\|f\|_2.
 \label{eq:packet-summed-tail}
\end{equation}
The spatial cutoff around \(X_J\) can be chosen larger by a fixed
factor than all the back-transported endpoint centers.  Outside it,
\(R\ge c\eps^{-\Delta}\), where
\(\Delta=u_r-b_*\).  The transport cutoff takes
\(R=\eps^{-\eta}\).  Thus their pairing errors are respectively
\(C_M\eps^{\Delta(M-1)}\|f\|_2\) and
\(C_M\eps^{\eta(M-1)}\|f\|_2\).
The same square-sum argument applies to a subset of coefficients.

There are at most \(C\eps^{-Q}\) retained coefficients.
Their part below \(\eps^B\|f\|_2^2\) has square sum at most
\(C\eps^{B-Q}\|f\|_2^2\).  Synthesis Bessel followed by the norm
bound for one test controls its pairing error by
\(C\eps^{(B-Q)/2}\|f\|_2\).
Finally, summing squared pairing errors over at most
\(C\eps^{-Q}\) tests costs a factor \(C\eps^{-Q}\).
The conservative choices \eqref{eq:packet-tail-orders} dominate
all these factors.  Increasing \(\eps_0^{-1}\) absorbs fixed constants
and proves all assertions.
\end{proof}

The tolerance in \eqref{eq:packet-transport} is \(\sigma_r\), the
\emph{earlier} position spacing.  At the endpoint uncertainty scale it
can be larger than \(\sigma_t\) by a power.  Thus this lemma alone
does not identify coherent packet masses with a positive law of
straight trajectories at that finer precision.

\subsection{Classes with cap count bounds}

Here are explicit conventions for the finite errors in the remaining
statements.  Let \(\xi>0\) be the mesh for coefficient and count
exponents, and let \(S\ge1\) bound the number of extra subdivision
and winner tags at each geometric localizer.  Constants implicit in
\(S\) may include powers of \(1+\ell\) from dyadic shape menus.
With fixed data, \(\eta,\xi,A\), an admissible packet error is
\begin{equation}
 e_{\rm pkt}(\eps)=
 C(\eta+\xi+\logeps S)
 +C_{\mathrm{data},\eta,\xi,A}
       \frac{1+\log(2+\ell)}{\ell}.
 \label{eq:packet-finite-error}
\end{equation}
The constants may be increased a finite number of times.  There is
no convergence assertion uniform as \(\eta\) tends to zero: the
order is first \(\eps\to0\), then \(\eta,\xi\to0\), with
\(\logeps S\to0\).  If a prescribed final error is needed, fix its
small \(\eta,\xi\) first and choose all integration orders and
\(\eps_0\) afterwards.

\begin{lemma}[Greedy cap classes]
\label{lem:cap-class}
In a geometric localizer, bin nonnegligible coefficients so that
\begin{equation}
 |\logeps{|c_T|^2}-\logeps m|\le\xi
 \label{eq:packet-mass-bin}
\end{equation}
on each bin.  Each bin admits a partition into classes, with a
count parameter \(g\) and fixed assigned-cap widths, having the
following properties.  If \(N_{\rm all}\) is the full class size,
then for every allowed local root cap \(F\),
\begin{equation}
 \logeps{\#\{T\hbox{ in the class}:T\in F\}}+w(F)
 \le g+e_{\rm pkt}.
 \label{eq:packet-cap-count}
\end{equation}
An empty count imposes no restriction.  There is an assigned cap
\(E=E(T)\) whose nonempty fibers in the full class satisfy
\begin{equation}
 \logeps{\#\{T:E(T)=E\}}
 \ge g-w(E)-e_{\rm pkt}.
 \label{eq:packet-assigned-fiber}
\end{equation}
There are only a number of classes polynomial in \(1+\ell\), with
constants depending on the fixed exponents and \(\xi\).
Any further tagged subdivision preserves
\eqref{eq:packet-cap-count}; its retained size is denoted by
\(N\le N_{\rm all}\).  The lower bound
\eqref{eq:packet-assigned-fiber} always refers to the full class.
\end{lemma}

\begin{proof}
Use the finite cap menu at local root widths up to
\(\min(L-r-b_*,u_r-b_*)\), with bounded enlargements so every
allowed cap is covered by boundedly many menu caps of the same
widths up to fixed factors.  On the current remaining population
\(\mathcal R\), choose a cap maximizing
\[
 \logeps{\#(\mathcal R\cap E)}+w(E).
\]
Assign its entire remaining contents to it and remove them.
The successive maxima
decrease.  Group removals by an exponent interval of length \(\xi\)
for this maximum and by the two assigned widths.  Call an upper
endpoint of the maximum interval \(g\).

The union in one such class lies inside the population present at
its first removal.  At that stage every cap count plus its weight
is at most \(g\), which proves \eqref{eq:packet-cap-count}.
For a removal belonging to this class, its own maximum is at least
\(g-\xi\), proving \eqref{eq:packet-assigned-fiber}.
Replacing an arbitrary cap by boundedly many menu caps changes the
normalized bound by at most a fixed constant divided by \(\ell\).
Rounding widths changes the weight by the same amount.  The
available coefficient exponents, after subtracting
\(\logeps{\|f\|_2^2}\), lie in a fixed bounded interval
after the cutoff in Lemma~\ref{lem:packet-decomposition}; so do
the maxima, since counts have exponent at most \(Q\) and weights
are bounded at fixed depths.  The number of bins is therefore
bounded in terms of \(\xi\), and the dyadic width choices contribute
only \(O((1+\ell)^2)\).  Finally, subdividing can only decrease
each cap count.
\end{proof}

For an endpoint observation in a localizer, choose a winning class
by the largest absolute pairing with its class superposition.
Here and later it is enough to keep endpoint masses satisfying
\begin{equation}
 m_O\ge\eps^{A_0}\|f\|_2^2
 \label{eq:packet-output-floor}
\end{equation}
for a sufficiently large fixed \(A_0\).  To see this without any
assumption on the law, normalize \(\|f\|_2=1\).  The number of
endpoint sites is at most \(C\eps^{-Q}\), and \(w\) is bounded
by a fixed constant \(W\).  A law on masses smaller than the
displayed threshold has score at most
\(\frac32 Q+\frac12W-\frac32A_0+O(1/\ell)\), whereas \(K_0\ge0\).
Choose \(A_0>Q+W+2\).  Splitting any law into this part and its
complement changes its averaged score by at most
\(3\log 2/(2\ell)\).  The low-mass part consequently cannot
increase a nonnegative upper growth bound.

Choose the tail order with \(A>A_0/2+1\).  Write
\(q=C_A\eps^{A-A_0/2}\), and reduce \(\eps_0\) so \(q<1/2\).
If there are at most \(S\) possible winning parts, the triangle
inequality gives the finite bound
\begin{equation}
 m_O(f)\le (1-q)^{-2}S^2m_O(h_J^{\rm win})
          \le4S^2m_O(h_J^{\rm win}).
 \label{eq:packet-winner-mass}
\end{equation}
Its loss in the mass term of the score is at most
\(3\logeps S+3\logeps{(1-q)^{-1}}\).
Thus no power-sized tail is discarded without comparison to an
output mass floor.  The label \(J\) below consists of \(J_0\),
the mass and cap bins, and any extra subdivision and winning tags.
The geometric part is determined by both endpoint data and aperture;
the additional tag entropy is bounded by \(\logeps S\).

\subsection{The local root and the earlier test}

\begin{proposition}[Local root comparison]
\label{prop:local-root-comparison}
Fix a winning retained class as above, and use translated coordinates
\eqref{eq:packet-rescale-coordinates}.  Let \(K_0^{\rm loc}\) be the
canonical root supremum of its input, including the norm option.
Set
\begin{equation}
 A_J=\logeps N+\frac32\logeps m+\frac12g.
 \label{eq:packet-local-root-value}
\end{equation}
Then, with the error convention \eqref{eq:packet-finite-error},
\begin{equation}
 K_0^{\rm loc}\le A_J+e_{\rm pkt}.
 \label{eq:packet-local-root-upper}
\end{equation}
For any chosen marginal law on the tagged localizers there is an
earlier test, using the original input at time \(r\), with score
\begin{equation}
 P_r^{\rm all}\ge\Ent(J)+b_*+
 \EE\left[\logeps{N_{\rm all}}+\frac32\logeps m+\frac12g\right]
 -e_{\rm pkt}.
 \label{eq:packet-earlier-whole-score}
\end{equation}
At \(r=0\), the right hand side is a lower bound for \(K_0\)
up to the displayed error.  The earlier test uses the uniform law
on each \emph{whole} class before the final retained subdivision.
\end{proposition}

\begin{proof}
The local root natural spacings are both
\(s=\eps^{u_r-b_*}\).  Let \(M_{\rm tot}\) be the sum of all
canonical root masses of the retained superposition, and let
\(M_F\) be their sum in a root cap \(F\).
By analysis and synthesis Bessel and the mass bin,
\begin{equation}
 M_{\rm tot}\le C\eps^{-\xi}Nm.
 \label{eq:packet-root-total-mass}
\end{equation}
Take \(R=\eps^{-\eta}\) in \eqref{eq:packet-cap-bessel}.
Its near set is covered by at most \(C\eps^{-4\eta}\) root
caps of the same widths.  Lemma~\ref{lem:cap-class} therefore
bounds its coefficient square sum by
\(C m\eps^{w(F)-g-\xi-4\eta-e_{\rm pkt}}\).
The far term is at most
\(C_M\eps^{\eta(2M-4)}\eps^{-\xi}Nm\).
Both \(N\) and the possible differences \(w(F)-g\) have bounded
exponents.  Choosing \(M\) sufficiently large, with those bounds
fixed, absorbs the far term in the same estimate, with an additional
fixed constant.  Consequently
\begin{equation}
 \logeps{M_F}\le\logeps m+g-w(F)+e_{\rm pkt}.
 \label{eq:packet-root-cap-mass}
\end{equation}
All canonical tests up to the logarithmic position cutoff are
covered by the uniform Bessel argument; no regularity constant grows
with this cutoff.

For an arbitrary root test law, the log-sum inequality gives
\[
 \Ent(O)+\EE\logeps{m_O}\le\logeps{M_{\rm tot}},\qquad
 \Ent(O\mid D)+\EE\logeps{m_O}\le\EE\logeps{M_D}.
\]
The second inequality is applied separately on each aperture label;
its observation support lies in that label's cap.  Adding the first
inequality and half the second, then adding \(\frac12\EE w(D)\),
gives \eqref{eq:packet-local-root-upper}.
The norm option is bounded by
\(\frac32(\logeps N+\logeps m)+e_{\rm pkt}\).
Zero-width caps cover the bounded local atom population in a fixed
number of pieces, so \eqref{eq:packet-cap-count} implies
\(\logeps N\le g+e_{\rm pkt}\).  This proves the same bound for
the norm option.

For the earlier test choose the prescribed marginal on \(J\),
then choose \(T\) uniformly in its whole class.  Take its earlier
site \(O_r=(C_r,V_T,X_T)\), and let \(D_r\) record \(J\) and
the assigned cap \(E(T)\).  Within \(J\), the map from atom
indices to earlier sites has bounded multiplicity.  Conversely,
an earlier site lies in only boundedly many geometric localizers,
because its center lies in a fixed enlargement of the cell in
\eqref{eq:packet-localizer}.  Extra tags cost at most \(\logeps S\).
It follows that
\[
 \Ent(O_r)\ge\Ent(J)+\EE\logeps{N_{\rm all}}-e_{\rm pkt},\qquad
 \Ent(O_r\mid D_r)\ge\EE[g-w(E)]-e_{\rm pkt}.
\]
The second inequality uses the full assigned fibers in
\eqref{eq:packet-assigned-fiber}, with their actual probabilities
under the uniform class law.  If one earlier site appears with
several coefficient tests, use the largest squared pairing there;
this is still an admissible test.  Thus its logarithmic mass is
at least the corresponding \(\logeps m-e_{\rm pkt}\).
At \(r=0\), \eqref{eq:packet-coeff-canonical} gives the same
conclusion with the canonical mass.  Finally, the earlier aperture
weight is \(2b_*+w(E)\), by the rescaling computation.  Substitution
in the score proves \eqref{eq:packet-earlier-whole-score}.
\end{proof}

\subsection{The preliminary growth bound}

\begin{proposition}[A finite upper bound for packet growth]
\label{prop:packet-trivial-growth}
The supremum \(\betaq\) is finite and \(\betaq\le1\).
This conclusion uses no identification of the original coherent
mass with a positive particle law.
\end{proposition}

\begin{proof}
Split into fixed aperture-width bins, at a cost tending to zero.
Apply Lemma~\ref{lem:packet-decomposition} with \(r=b_*=0\), then
the class and winner construction.  Restrict to the mass floor as
justified above.  Work within one tagged localizer and its retained
class.  By \eqref{eq:packet-winner-mass}, the original conditional
score is at most the score of the winning superposition plus
\(e_{\rm pkt}\).  For the remainder of this conditional calculation,
\(m_O\) denotes the winning-superposition masses.  They satisfy
the inherited floor
\[
 m_O\ge (1-q)^2S^{-2}\eps^{A_0}\|f\|_2^2.
\]
The factor \(3\logeps{(1-q)^{-1}}\) is included in
\(e_{\rm pkt}\) by reducing \(\eps_0\); for subpower tag counts
the displayed floor still has a fixed polynomial exponent.
Let \(Z=(C_t,P)\), where \(P\) records both velocity and
back-transported position \(X+(C_0-C_t)V\) isotropically at
depth \(a\).  Set \(m_Z=\sum_{O:O\mapsto Z}m_O\).
The log-sum inequality on each fiber shows
\begin{equation}
 \Ent(O)+\tfrac12\Ent(O\mid D)+\tfrac32\EE\logeps{m_O}
 \le \Ent(Z)+\tfrac12\Ent(Z\mid D)+\tfrac32\EE\logeps{m_Z}.
 \label{eq:packet-coarsen-mass}
\end{equation}
For clarity, the unconditional log-sum inequality bounds
\(\Ent(O\mid Z)+\EE\logeps{m_O/m_Z}\) by zero; the
same conditional on \((Z,D)\), using a sum over a subset of the
fiber, bounds its conditional counterpart by zero.  Their sum with
coefficients one and one half is precisely the displayed inequality.

Let \(\mathcal T_Z\) consist of all class atoms compatible with
\(P\) up to the enlargement in \eqref{eq:packet-transport}.
The earlier position uncertainty and the endpoint velocity
uncertainty change these coarse columns by at most
\(C\eps^{-\eta}\eps^a\), since \(a\le u_t\le u_0\).
Discarding interactions outside \(\mathcal T_Z\) is permitted by
the quantitative tail lemma.  After fixing the coefficient cutoff
\(B\), increase only the transport integration order so its error
exponent \(A_{\rm tr}\) also permits summation over this fiber and
comparison to that coefficient floor.  This does not require
another coefficient truncation.  Bessel at the fixed time \(C_t\)
then gives
\begin{equation}
 m_Z\le C\eps^{-\xi}m\,\#\mathcal T_Z
 \label{eq:packet-contributor-mass}
\end{equation}
after absorbing those tails.  More explicitly, before absorption
the right hand side has an added term
\(C_{A_{\rm tr}}\eps^{2A_{\rm tr}-Q}\|f\|_2^2\); choosing
\(2A_{\rm tr}-Q>B+1\), and then reducing \(\eps_0\), absorbs it whenever
\(\mathcal T_Z\ne\varnothing\).  The output floor excludes a
nonempty retained fiber with no potential contributor.  The
constant in \eqref{eq:packet-contributor-mass} also absorbs the
fixed factor from squaring the sum of a main term and its error.

Only now introduce a positive auxiliary law: keep the given law
of \((Z,D)\), and choose \(T\) uniformly in \(\mathcal T_Z\),
independently of \(D\) conditional on \(Z\).  Exactly,
\[
 \Ent(T\mid Z,D)=\Ent(T\mid Z)
  =\EE\logeps{\#\mathcal T_Z}.
\]
The chain rule gives the contributor calculation
\begin{align}
 &\Ent(Z)+\tfrac12\Ent(Z\mid D)
                      +\tfrac32\Ent(T\mid Z)\nonumber\\
 &=\Ent(T)+\tfrac12\Ent(T\mid D)
                  +\Ent(Z\mid T)+\tfrac12\Ent(Z\mid T,D)\nonumber\\
 &\le\Ent(T)+\tfrac12\Ent(T\mid D)+\Ent(C_t\mid T)
                     +\tfrac32\Ent(Z\mid T,C_t).
 \label{eq:packet-contributor-identity}
\end{align}
Here \(C_t\) is determined by both \(Z\) and \(D\).  Thus
\(\Ent(Z\mid T)=\Ent(C_t\mid T)+\Ent(Z\mid T,C_t)\), and
\(\Ent(Z\mid T,D)\le\Ent(Z\mid T,C_t)\).
This keeps exactly one conditional time entropy.  The remaining
term measures the uncertainty in the coarse observation after the
contributor and time are known.

In the present root-referred observation, \((T,C_t)\) locates \(Z\)
to at most \(C\eps^{-C\eta}\) possibilities.  Hence
\(\Ent(Z\mid T,C_t)\le C\eta+O(1/\ell)\).
Combining \eqref{eq:packet-contributor-identity} with
\eqref{eq:packet-contributor-mass} yields
\begin{equation}
 \Ent(Z)+\tfrac12\Ent(Z\mid D)+\tfrac32\EE\logeps{m_Z}
 \le\Ent(T)+\tfrac32\logeps m+
       \tfrac12\Ent(T\mid D)+\Ent(C_t\mid T)+e_{\rm pkt}.
 \label{eq:packet-trivial-identity}
\end{equation}

For a fixed aperture label \(D\), its potential contributors lie
in a \(\eps^{-C\eta}\) enlargement of the root aperture of widths
\((b,a)\) obtained by referring its rectangles back to \(C_0\).
The velocity error is at most \(C\rho_t\); the position error is
at most \(C\eps^{-\eta}\sigma_0+C\rho_t\).  Both are bounded
by the corresponding enlarged coarse widths, because \(a\le u_t\).
Covering this enlargement and applying
\eqref{eq:packet-cap-count} gives
\[
 \Ent(T\mid D)\le g-\EE w(D)+e_{\rm pkt},\qquad
 \Ent(T)\le\logeps N,\qquad
 \Ent(C_t\mid T)\le t+O(1/\ell).
\]
Adding the half-weight to \eqref{eq:packet-trivial-identity}
therefore bounds the conditional endpoint score by
\(A_J+t+e_{\rm pkt}\).  Restoring \(J\) adds
\(\Ent(J)\), with the recorded tag error; here \(b_*=0\).
Proposition~\ref{prop:local-root-comparison} at \(r=0\) gives
\[
 P_t\le\Ent(J)+\EE A_J+t+e_{\rm pkt}
 \le K_0+t+2e_{\rm pkt},
\]
because \(N\le N_{\rm all}\).  First take \(\eps\to0\) at fixed
depths, \(\eta,\xi\) and integration orders, then send
\(\eta,\xi\to0\).  This proves \(\betaq\le1\).
\end{proof}
